\chapter{正规族}

\begin{solution}{normal-local-uniform-integrals}
设 $\gamma:[a,b]\to\Omega$ 分段 $C^1$。其像 $K=\gamma([a,b])$ 是
$\Omega$ 中的紧集，故 $f_n\to f$ 在 $K$ 上一致。于是
\[
 \left|\int_\gamma f_n(z)\,dz-\int_\gamma f(z)\,dz\right|
 \le \sup_K|f_n-f|\int_a^b|\gamma'(t)|dt\longrightarrow0.
\]
分段点不影响长度估计。
\end{solution}

\begin{solution}{normal-derivative-convergence}
设 $K\Subset\Omega$，取 $\rho>0$ 使闭 $\rho$ 邻域
$K_\rho\Subset\Omega$。对 $z\in K$ 在圆 $|\zeta-z|=\rho$ 上使用 Cauchy
导数公式，得
\[
 |f_n'(z)-f'(z)|
 \le \frac1\rho\sup_{\zeta\in K_\rho}|f_n(\zeta)-f(\zeta)|.
\]
右端与 $z$ 无关并趋于零。高阶导数同理，常数变为 $k!\rho^{-k}$。
\end{solution}

\begin{solution}{normal-power-sequence}
若 $K\Subset\mathbb D$，则 $r:=\max_K|z|<1$，所以
$\sup_K|z^n|\le r^n\to0$。另一方面，
\[
 \sup_{|z|<1}|z|^n=1
\]
对每个 $n$ 都成立，故在整个开圆盘上的上确界误差不趋零。注意上确界不必
在开圆盘内取得。
\end{solution}

\begin{solution}{normal-compact-open-metric}
先核验度量公理。级数被 $\sum2^{-m}=1$ 控制，非负且对称。
距离为零当且仅当各紧集上函数相等，因穷竭覆盖区域，这等价于两函数相等。
对每个 $m$ 使用上确界范数三角不等式，再用
\[
 \min(1,s+t)\le\min(1,s)+\min(1,t)\qquad(s,t\ge0),
\]
乘以 $2^{-m}$ 并求和，得到 $d(f,h)\le d(f,g)+d(g,h)$。

若 $f_n\to f$ 局部一致，则每个固定 $m$ 的
$\|f_n-f\|_{C(K_m)}\to0$。给定 $\varepsilon>0$，先取 $M$ 使
$\sum_{m>M}2^{-m}<\varepsilon/2$，再同时控制前 $M$ 项，得到
$d(f_n,f)<\varepsilon$。

反之，若 $d(f_n,f)\to0$，对固定 $m$ 有
\[
 2^{-m}\min\{1,\|f_n-f\|_{C(K_m)}\}\le d(f_n,f)\to0,
\]
故在 $K_m$ 上一致收敛。任意紧集包含于某个 $K_m$ 的内部，因而得到局部
一致收敛。两套穷竭都刻画“在每个紧集上一致收敛”，所以所得收敛概念相同。
\end{solution}

\begin{solution}{normal-cauchy-equicontinuity}
取 $K\Subset U\Subset\Omega$，并令
$\delta=\operatorname{dist}(K,\partial U)>0$。在 $U$ 上统一有 $|f|\le M$。
以半径 $\delta/4$ 的圆盘有限覆盖 $K$。若 $z,w$ 位于同一个这样的圆盘，
连接线段落在以同一圆心、半径 $\delta/2$ 的圆盘中；Cauchy 估计在那里给
$|f'|\le 2M/\delta$。因此
\[
 |f(z)-f(w)|\le(2M/\delta)|z-w|.
\]
若两点不在同一覆盖圆盘，而 $|z-w|$ 足够小，可先取覆盖的 Lebesgue 数，
保证二点集落入某个覆盖成员。所得模数与 $f$ 无关，故等度连续。
\end{solution}

\begin{solution}{normal-disk-maps}
因 $|f|<1$，函数族局部有界，Montel 定理给正规性。设子列局部一致收敛到
$g$。则 $|g|\le1$。若存在 $a\in\Omega$ 使 $|g(a)|=1$，最大模原理说明
$g$ 为常数，且该常数位于单位圆周。否则 $|g(z)|<1$ 对所有 $z$ 成立，故
$g:\Omega\to\mathbb D$。所以极限或者仍映入开圆盘，或者退化为单位圆周上
的常值函数。
\end{solution}

\begin{solution}{normal-basepoint-derivatives}
对非空 $K\Subset\Omega$，选有限个圆盘 $B(z_j,r_j)$ 覆盖 $K$，
使各闭圆盘包含于 $\Omega$。平面域可由折线连通，故有从 $a$ 到各 $z_j$
的域内折线 $\alpha_j$。令 $L$ 为这些折线与闭圆盘的有限并，它是
$\Omega$ 中的紧集。若 $z\in B(z_j,r_j)$，把 $\alpha_j$ 接上从 $z_j$
到 $z$ 的直线段，得到全在 $L$ 内、长度不超过
$C_L:=\max_j(\ell(\alpha_j)+r_j)$ 的路径。
这是由具体构造得到的长度界，不能仅从紧致道路连通性推出。
导函数族在 $L$ 上一致有界，故
\[
 |f(z)|\le |f(a)|+\left|\int_a^z f'(\zeta)d\zeta\right|
 \le C_0+C_LM_L.
\]
这给出 $K$ 上一致界，不需要 $\Omega$ 凸。
\end{solution}

\begin{solution}{normal-vitali-dense-set}
稠密子集只要非空，就在区域内有聚点；更直接地，区域中任意小圆盘含有无穷
多个稠密集的点。题设因此满足 Vitali--Porter 定理的集合条件，局部有界性
给出所需正规性，故整个序列局部一致收敛到全纯函数。
\end{solution}

\begin{solution}{normal-vitali-necessity}
每个 $p_n$ 是多项式。若 $z=0$，因前置因子 $z$，恒有 $p_n(0)=0$；若
$z=1/k$，则当 $n\ge k$ 时乘积中出现因子 $z-1/k=0$。所以在
$E=\{0\}\cup\{1/k:k\ge2\}\subset\mathbb D$ 上逐点极限为零，且 $0$ 是
位于单位圆盘内的聚点。

但在圆盘内部的 $z=-1/2$，
\[
 |p_n(-1/2)|
 =\frac{n!}{2}\prod_{k=1}^n\left(\frac12+\frac1k\right)
 =\frac{n!(n+1)(n+2)}{2^{n+2}}\to\infty.
\]
因此函数列在含 $-1/2$ 的任何紧邻域上不一致有界，也不可能局部一致收敛。
\end{solution}

\begin{solution}{normal-finite-normal-converse}
反设某紧集 $K\Subset\Omega$ 上不存在统一界。可递归选取
$f_n\in\mathcal F$ 与 $z_n\in K$，使 $|f_n(z_n)|>n$。由紧性，取子列使
$z_{n_j}\to z\in K$。正规性再给进一步子列 $f_{n_{j_\ell}}$ 局部一致收敛
到有限全纯函数 $f$。取紧邻域 $L\Subset\Omega$ 包含 $z$ 和所有充分大的
$z_{n_{j_\ell}}$，局部一致收敛说明这些函数在 $L$ 上一致有界，矛盾。
\end{solution}

\begin{solution}{normal-holomorphic-series}
记部分和 $S_N=\sum_{n=0}^Nf_n$。题设说 $(S_N)$ 局部有界，并在一个具有
内部聚点的集合上逐点收敛。Vitali--Porter 定理给 $S_N\to S$ 局部一致，
且 $S$ 全纯。由全纯极限定理，
\[
 S_N'=\sum_{n=0}^N f_n'\longrightarrow S'
\]
局部一致。因此原级数可逐项求导，导数级数也在紧集上一致收敛。
\end{solution}

\begin{solution}{normal-coefficient-compactness}
函数族映入单位圆盘，故由 Montel 定理可取子列 $f_{n_j}$ 局部一致收敛到
$f\in\mathcal O(\mathbb D)$。全纯极限定理给每阶导数在零点收敛，因而
\[
 a_{n_j,k}=\frac{f_{n_j}^{(k)}(0)}{k!}
 \longrightarrow \frac{f^{(k)}(0)}{k!}=:a_k
\]
对每个固定 $k$ 成立。所得 $(a_k)$ 正是 $f$ 在零点的 Taylor 系数；因此它
不是逐坐标抽取产生的任意形式序列，而确实在单位圆盘内表示全纯极限
$f(z)=\sum a_kz^k$。
\end{solution}

\begin{solution}{normal-univalent-closure}
局部一致极限 $f$ 全纯，且导数也局部一致收敛，所以
\[
 |f'(a)|=\lim_n|f_n'(a)|\ge c>0.
\]
故 $f$ 不是常函数。单叶函数极限的推论随即说明 $f$ 仍单叶。若没有导数
下界，极限可能退化为常值，例如 $f_n(z)=z/n$ 在单位圆盘上均单叶，却局部
一致收敛到零。
\end{solution}
